\chapter{État de l'art}
\label{chap:etat-art}

\section{Introduction}
\label{sec:ea-intro}

Avant de concevoir une architecture de sécurisation des flux IoT, il est indispensable de cadrer le sujet sur trois plans complémentaires~: (i)~caractériser les \textbf{menaces} qui pèsent spécifiquement sur les écosystèmes IoT à grande échelle~; (ii)~recenser et évaluer de manière critique les \textbf{technologies et protocoles} candidats pour le transport, l'authentification et la supervision~; (iii)~analyser les \textbf{travaux et solutions existants} afin d'en extraire les bonnes pratiques, d'identifier leurs limites et de positionner notre contribution. Ce chapitre est structuré autour de ces trois axes, suivi d'une analyse critique de l'existant et d'une synthèse positionnant le projet.

\section{Menaces et vulnérabilités spécifiques à l'IoT}
\label{sec:ea-menaces}

\subsection{Surface d'attaque étendue}

Les écosystèmes IoT présentent une surface d'attaque significativement plus large que les systèmes informatiques classiques. Meneghello \emph{et al.}~\cite{meneghello2019iot} en distinguent quatre couches~: (i)~la couche \emph{matérielle} (composants embarqués, JTAG/UART exposés, attaques par canal auxiliaire) . (ii)~la couche \emph{firmware} (mises à jour non signées, exécution de code arbitraire) . (iii)~la couche \emph{réseau} (transports en clair, protocoles propriétaires faibles) . (iv)~la couche \emph{services et écosystème} (\acs{API} non authentifiées, applications mobiles vulnérables, multi-tenance mal cloisonnée).

\subsection{OWASP IoT Top 10 et menaces réseau}

L'OWASP synthétise dans son \emph{IoT Top~10}~\cite{owasp_iot} les dix risques applicatifs majeurs~: mots de passe faibles ou par défaut, services réseau non sécurisés, interfaces écosystème non protégées, mécanisme de mise à jour absent, composants obsolètes, protection insuffisante de la vie privée, transferts et stockage non chiffrés, absence de gestion des dispositifs, paramètres par défaut non sûrs, absence de hardening physique. Dans ce projet, nous nous concentrons sur les menaces \textbf{réseau} et \textbf{applicatives}~: interception (\emph{sniffing}), \emph{man-in-the-middle}, force brute, exfiltration applicative, déni de service, vol/abus de certificats.

\subsection{Études de cas de référence}

L'analyse de l'attaque \textbf{Mirai}~\cite{kolias2017mirai,antonakakis2017mirai} illustre toutes ces faiblesses simultanément~: le botnet exploitait des couples identifiant/mot de passe par défaut sur des objets exposés en SSH/Telnet, recrutait des centaines de milliers de dispositifs en quelques jours, et générait des attaques DDoS dépassant 1~Tbit/s. Le rapport \emph{ENISA Threat Landscape for IoT}~\cite{enisa_iot} confirme que ces vecteurs restent dominants et identifie comme contre-mesures prioritaires~: l'authentification forte, le chiffrement systématique, la segmentation réseau et la supervision active. Ces orientations rejoignent les recommandations du NIST~\cite{nist_sp800213} pour la cybersécurité des dispositifs IoT.

Cette cartographie des menaces permet d'identifier les vecteurs prioritaires à contrer dans la solution~: interception des flux, usurpation d'identité, déni de service et abus de certificats. Le premier levier sur lequel il est possible d'agir est le \textbf{protocole de communication applicatif} lui-même~: c'est à ce niveau que la surface d'attaque réseau est déterminée et que les protections de base peuvent être imposées.

\section{Protocoles applicatifs IoT}
\label{sec:ea-protocoles}

Le tableau~\ref{tab:protocoles} compare quatre protocoles applicatifs représentatifs.

\begin{table}[!htbp]
\centering
\caption{Comparaison des principaux protocoles applicatifs IoT}
\label{tab:protocoles}
\small
\begin{tabularx}{\linewidth}{|l|X|X|X|X|}
\hline
\textbf{Critère} & \textbf{MQTT 5} & \textbf{CoAP} & \textbf{AMQP 1.0} & \textbf{HTTP/REST} \\
\hline
Modèle & Pub/Sub & Req/Rep & Pub/Sub + queues & Req/Rep \\ \hline
Transport & TCP/TLS & UDP/DTLS & TCP/TLS & TCP/TLS \\ \hline
Empreinte & Très faible & Très faible & Moyenne & Moyenne/élevée \\ \hline
QoS & 0/1/2 & Confirmable & 0/1 & Aucun natif \\ \hline
Sécurité native & TLS+ACL & DTLS+OSCORE & SASL+TLS & TLS+JWT \\ \hline
Maturité IoT & Très large & Large & Industrielle & Universelle \\ \hline
\end{tabularx}
\end{table}

Le choix de \textbf{MQTT~5}~\cite{mqtt5} pour ce projet repose sur trois critères~: empreinte mémoire et réseau adaptée aux microcontrôleurs . large adoption dans les écosystèmes industriels et les opérateurs IoT . existence de mécanismes natifs d'authentification (\acs{TLS}+\acs{ACL}) facilement combinables avec mTLS.

Le protocole applicatif étant fixé, la question du \textbf{chiffrement du canal} devient prépondérante~: MQTT, en lui-même, ne garantit ni la confidentialité ni l'authenticité des messages échangés. TLS~1.3 avec authentification mutuelle répond précisément à ces deux exigences, comme le détaille la section suivante.

\section{Sécurisation du transport : TLS et mTLS}
\label{sec:ea-tls}

\subsection{TLS 1.3}

Le protocole \acs{TLS}~1.3~\cite{rfc8446} apporte par rapport à TLS~1.2 une simplification significative~: suppression des suites cryptographiques obsolètes, négociation en \textbf{1-RTT} (et 0-RTT optionnel), \emph{forward secrecy} obligatoire, signatures et échanges de clés modernisés (ECDHE, EdDSA, AES-GCM, ChaCha20-Poly1305). Les recommandations opérationnelles publiées dans la \acs{RFC}~9325~\cite{rfc9325} guident le choix des paramètres et la durée de vie des certificats.

\subsection{Authentification mutuelle}

Dans un déploiement classique, seul le serveur s'authentifie auprès du client. Le \textbf{mTLS} (\emph{mutual TLS}) impose en outre au client de présenter un certificat X.509 valide signé par une autorité de confiance. Cette double authentification permet~: (i) d'éviter qu'un objet illégitime se connecte au broker . (ii) de chiffrer chaque session par bout . (iii) de tracer chaque message via l'identité cryptographique de son émetteur. Cette propriété est essentielle pour la mise en place ultérieure d'\acs{ACL} par certificat dans Mosquitto~\cite{mosquitto_tls}.

L'authentification mutuelle par certificats suppose cependant l'existence d'une infrastructure capable d'\textbf{émettre, renouveler et révoquer} ces certificats à l'échelle d'un parc de dispositifs~: c'est précisément le rôle d'une PKI automatisée, dont les principes sont analysés dans la section suivante.

\begin{resultbox}{Positionnement critique~: TLS 1.3 + mTLS}
Le choix de TLS~1.3 en remplacement de TLS~1.2 répond à trois arguments techniques précis~: (i)~la suppression des suites vulnérables (RC4, DES, RSA pur sans \emph{forward secrecy}) réduit la surface cryptographique exploitable~; (ii)~la négociation en 1-RTT diminue la latence du \emph{handshake}, critique pour les dispositifs contraints~; (iii)~la \emph{forward secrecy} obligatoire garantit que la compromission ultérieure d'une clé privée ne met pas en péril les sessions passées. L'ajout du mTLS va au-delà du chiffrement~: il substitue à une authentification par mot de passe (facilement rejouable) une preuve cryptographique d'identité non duplicable sans accès à la clé privée du dispositif. Cette double garantie~-- confidentialité \emph{et} authenticité mutuelle~-- est la seule combinaison capable de fermer simultanément les vecteurs d'interception et d'usurpation identifiés dans la cartographie des menaces.
\end{resultbox}

\section{Gestion des identités : PKI et automatisation}
\label{sec:ea-pki}

Une PKI à deux niveaux (CA racine \emph{offline} + CA intermédiaire \emph{online}) constitue la pratique recommandée~\cite{nist_sp800213}~: la CA racine ne signe que des CA intermédiaires, et reste hors ligne . les CA intermédiaires signent les certificats opérationnels (services, dispositifs). En cas de compromission d'une CA intermédiaire, seul le sous-arbre concerné doit être révoqué.

L'automatisation de l'émission est indispensable lorsque le parc dépasse quelques dizaines d'objets. \textbf{HashiCorp Vault}~\cite{vault_pki} fournit un \emph{secrets engine} \texttt{pki} qui expose une API HTTP de demande de certificats avec contrôle fin (rôles, ACL, TTL). C'est la solution retenue dans ce projet.

\begin{resultbox}{Positionnement critique~: Vault vs alternatives PKI}
Trois alternatives à HashiCorp Vault ont été évaluées~: (i)~\textbf{OpenSSL CLI}~: complet techniquement mais entièrement manuel~; toute émission de certificat requiert une intervention humaine, ce qui l'exclut dès que le parc dépasse une dizaine de dispositifs~; (ii)~\textbf{EJBCA (Keyfactor)}~: solution PKI d'entreprise mature, mais sa complexité de déploiement et sa courbe d'apprentissage sont disproportionnées pour un projet académique de six mois~; (iii)~\textbf{Let's Encrypt / ACME}~: idéal pour les certificats TLS publics, mais inadapté à la PKI interne d'un réseau IoT isolé (absence de validation domaine interne, pas de certificats client mTLS natifs). Vault a été retenu car il est le seul à réunir~: API REST d'émission automatisée, gestion fine des rôles et TTL, intégration aux pipelines DevOps, et déploiement entièrement on-premise~-- critère non négociable dans un contexte opérateur télécom soucieux de souveraineté des données.
\end{resultbox}

Même avec une PKI robuste et un chiffrement bout en bout, il n'est pas possible d'exclure tout incident~: un certificat peut être volé, un dispositif légitime peut adopter un comportement déviant, ou une configuration insuffisante peut ouvrir un vecteur non anticipé. La \textbf{supervision active}~-- combinant IDS réseau et SIEM~-- constitue le filet de sécurité complémentaire à ces mécanismes préventifs.

\section{Détection d'intrusion et supervision}
\label{sec:ea-ids-siem}

\subsection{IDS réseau orienté MQTT}

\textbf{Suricata} est un IDS/IPS open source qui dispose, depuis la version 6, d'un \emph{parser} natif MQTT capable d'extraire les champs \texttt{CONNECT}, \texttt{PUBLISH}, \texttt{SUBSCRIBE}~\cite{suricata_mqtt}. Cette capacité permet de définir des règles métier (par exemple~: alerter sur un \texttt{CONNECT} sans \texttt{ClientID} reconnu, ou sur une rafale de \texttt{CONNECT} indiquant un \emph{brute force}) qui ne seraient pas accessibles à un IDS générique.

Le tableau~\ref{tab:ids-compare} compare trois solutions IDS open source représentatives sur les critères déterminants pour ce projet.

\begin{table}[!htbp]
\centering
\caption{Comparaison des solutions IDS open source}
\label{tab:ids-compare}
\small
\begin{tabularx}{\linewidth}{|l|X|X|X|}
\hline
\textbf{Critère} & \textbf{Suricata 6} & \textbf{Snort 3} & \textbf{Zeek 6} \\
\hline
Parser MQTT natif    & Oui (v6+)     & Non           & Partiel (scripts) \\ \hline
IDS/IPS actif        & Oui           & Oui           & Non (analyse seule) \\ \hline
Multi-threading      & Oui (natif)   & Oui (v3)      & Oui \\ \hline
Intégration Wazuh    & Native (eve.json) & Adaptateur requis & Scripts Lua \\ \hline
Règles Sigma/Suricata & Native       & Conversion requise & Non \\ \hline
Maturité IoT         & Élevée        & Moyenne       & Élevée (analytique) \\ \hline
\textbf{Verdict}     & \textbf{Retenu} & Non retenu   & Complémentaire \\ \hline
\end{tabularx}
\end{table}

Suricata est retenu pour trois raisons décisives~: (i)~son parser MQTT natif élimine tout besoin de post-traitement applicatif~; (ii)~son format de sortie \texttt{eve.json} est nativement consommé par Wazuh, réduisant la complexité d'intégration~; (iii)~son moteur multi-thread garantit une analyse sans dégradation de débit sur les segments IoT à forte densité de messages.

\subsection{SIEM Wazuh}

\textbf{Wazuh}~\cite{wazuh_docs} est une plateforme de \emph{Security Information and Event Management} qui agrège les journaux d'agents endpoint, de sources réseau (Suricata) et applicatifs (Vault, Mosquitto). Elle s'intègre nativement à OpenSearch pour le stockage et au tableau de bord Wazuh Dashboard pour la visualisation et la corrélation. Sa modularité (règles XML, décodeurs personnalisables) en fait un candidat naturel pour la consolidation des alertes IoT.

Le tableau~\ref{tab:siem-compare} compare Wazuh aux deux autres plates-formes SIEM courantes sur les critères déterminants pour un projet IoT académique à contraintes de coût et de souveraineté.

\begin{table}[!htbp]
\centering
\caption{Comparaison des solutions SIEM}
\label{tab:siem-compare}
\small
\begin{tabularx}{\linewidth}{|l|X|X|X|}
\hline
\textbf{Critère} & \textbf{Wazuh 4.7} & \textbf{ELK Stack} & \textbf{Splunk Free} \\
\hline
Licence              & Open source (OSS)   & Open source    & Propriétaire (limité) \\ \hline
Intégration Suricata & Native (eve.json)   & Via Filebeat    & Via add-on HEC \\ \hline
Corrélation native   & Oui (règles XML)    & Non (Elastic SIEM en payant) & Oui (SPL) \\ \hline
Agents endpoint      & Oui (Wazuh agent)   & Non natif       & Oui (UF) \\ \hline
Volume journaux (gratuit) & Illimité       & Illimité        & 500 MB/j \\ \hline
Courbe d'apprentissage & Modérée          & Élevée (tuning) & Modérée \\ \hline
Souveraineté données & Totale (on-premise) & Totale          & Partielle (cloud) \\ \hline
\textbf{Verdict}     & \textbf{Retenu}     & Non retenu      & Non retenu \\ \hline
\end{tabularx}
\end{table}

Wazuh est retenu pour sa triple intégration native~: avec Suricata (format \texttt{eve.json} sans transformation), avec les agents endpoint sur les conteneurs Docker du laboratoire, et avec OpenSearch pour un stockage illimité. La licence open source garantit en outre la souveraineté des données, critère non négociable dans le contexte opérateur télécom.

Les règles IDS, aussi précises soient-elles, ne peuvent couvrir que les attaques \textbf{déjà connues et formalisées}. Pour détecter des comportements nouveaux ou subtils~-- comme l'utilisation d'un certificat valide depuis un contexte comportemental anormal~-- une couche d'apprentissage automatique est nécessaire. La section suivante en présente les approches retenues.

\section{Apprentissage automatique pour la détection d'anomalies}
\label{sec:ea-ml}

Les approches par règles statiques sont efficaces pour les attaques connues mais peinent à détecter les comportements nouveaux. Hasan \emph{et al.}~\cite{hasan2020securing} ont montré qu'un classifieur supervisé entraîné sur des features réseau et applicatives obtient des taux de détection supérieurs à 99~\% sur sept catégories d'attaques IoT. Les algorithmes les plus utilisés sont~:

\begin{itemize}[leftmargin=*]
  \item \textbf{Isolation Forest}~\cite{liu2008isolation}~: méthode non supervisée d'isolement, particulièrement efficace en haute dimension et adaptée aux scénarios où les anomalies sont rares .
  \item \textbf{Random Forest}~\cite{breiman2001random}~: ensemble d'arbres de décision, robuste, interprétable (\emph{feature importance}) et performant en classification multi-classes.
\end{itemize}

La mobilisation conjointe de ces deux algorithmes n'est pas arbitraire~: elle répond à deux exigences analytiquement distinctes. Isolation Forest est \textbf{non supervisé}~: il modélise la normalité sans nécessiter de données étiquetées, ce qui le rend indispensable pour détecter les anomalies \emph{inconnues}~-- typiquement, l'utilisation d'un certificat valide depuis un contexte comportemental inhabituel (scénario~S6). Random Forest est \textbf{supervisé}~: entraîné sur des exemples d'attaques étiquetées, il fournit une classification précise en sept catégories et une mesure d'importance par feature permettant d'identifier les attributs les plus discriminants. Ces deux algorithmes sont donc complémentaires~: l'un couvre l'inconnu, l'autre maximise la précision sur le connu. La bibliothèque \emph{scikit-learn}~\cite{scikit_learn} fournit des implémentations matures de ces deux algorithmes . elle constitue le socle de notre pipeline ML.

\section{Analyse critique de l'existant}
\label{sec:ea-analyse}

Plusieurs approches ont été proposées pour sécuriser les flux IoT~: gateways IoT propriétaires (AWS IoT Core, Azure IoT Hub), distributions intégrées (ThingsBoard, Eclipse Kura), ou architectures sur mesure. Notre analyse fait ressortir quatre limitations récurrentes~:

\begin{enumerate}[leftmargin=*]
  \item \textbf{Verrouillage fournisseur} pour les solutions cloud commerciales, peu compatible avec une stratégie souveraine d'opérateur télécom .
  \item \textbf{Couverture partielle} de la sécurité chez les distributions intégrées (souvent TLS oui, mais pas mTLS systématique ni PKI automatisée) .
  \item \textbf{Faible intégration} entre la couche réseau (IDS) et la couche applicative (broker, ML) .
  \item \textbf{Reproductibilité limitée} des architectures publiées dans la littérature, rarement accompagnées d'un environnement de simulation complet.
\end{enumerate}

\section{Synthèse du chapitre}
\label{sec:ea-positionnement}

Au terme de cet état de l'art, le présent projet se positionne non comme une juxtaposition d'outils, mais comme une \textbf{architecture de référence open source} pour la sécurisation des flux IoT dans un contexte opérateur. Quatre apports structurants s'en dégagent~: (1)~un déploiement de bout en bout reproductible sous GNS3 et Docker~; (2)~une PKI Vault automatisée associée à un usage systématique du mTLS~; (3)~une chaîne de détection unifiée Suricata~+~Wazuh enrichie de règles MQTT spécifiques~; (4)~un module ML complémentaire, combinant IsolationForest et RandomForest, pour traiter les anomalies peu ou pas couvertes par les signatures.

\section{Conclusion}
\label{sec:ea-conclusion}

Ce chapitre a établi le socle conceptuel du projet en suivant un fil directeur continu~: des menaces spécifiques à l'IoT découlent les exigences de protocole~; du choix de MQTT découle le besoin de TLS~1.3~; de TLS avec mTLS découle la nécessité d'une PKI automatisée~; de la PKI découle le besoin de supervision IDS/SIEM~; enfin, des limites des règles statiques découle l'apport de l'apprentissage automatique. Son rôle était de dégager les exigences techniques et les principes de protection, non de diagnostiquer encore la plateforme cible. Le chapitre suivant traduit ces constats théoriques en besoins opérationnels, puis en choix de conception adaptés au contexte étudié.
